\begin{exercise}{109}
Show that a Kan fibration $p\colon X\to Y$ which is additionally a weak
equivalence has contractible fibres.

\begin{proof}
Fix $y\in Y_0$ and put $F=X_y$, a Kan complex by pullback.
By Exercise~\ref{ex:108}, $q=|p|$ is a Serre fibration with fibre
$|F|$ over $|y|$.
The bijection on path components induced by $q$, together with path
lifting, makes this fibre nonempty.
Choose a basepoint in it.

The homotopy exact sequence of $q$ shows that
$\pi_r(|F|)=0$ for every $r\ge1$, since $q$ induces isomorphisms
on all homotopy groups.
For the component statement, injectivity of $\pi_0(q)$ implies that
all components of $|F|$ lie in the same component of $|X|$.
Thus the monodromy action of $\pi_1(|Y|,|y|)$ on $\pi_0(|F|)$ is
transitive. The stabilizer of the chosen component is the image
of $\pi_1(|X|)$, which is all of $\pi_1(|Y|)$.
Hence $|F|$ has one component.
Since $|F|$ is a CW complex (Corollary~1.1.41), Whitehead's theorem
now makes it contractible.

Applying $\mathcal S$ to a topological contraction and using the
canonical map $\Delta^1\to\mathcal S([0,1])$ gives a simplicial
contraction of $\mathcal S(|F|)$.
The Kan-unit fact granted in Exercise~\ref{ex:97} says that
$F\to\mathcal S(|F|)$ is a homotopy equivalence.
Transporting the contraction across this homotopy equivalence shows
that $F$ is contractible.
If $Y=\varnothing$, there are no fibres to check.
\end{proof}
\end{exercise}
